% !PRINTED 438
\weberpage{438}

%<<CONT
is therefore expressible \gesperrt{rationally} through $(k)$ and $(lk)$, and from § 72
it follows that $M$ is an \gesperrt{integral algebraic number}.

We apply this theorem to each pair of consecutive
terms of the sequence (7) and obtain, if we set % [sic] The German print cites sequence (7); it means the sequence numbered (10) on the preceding page.
\[
\omega_1 = \frac{\omega}{p}, \quad \omega_2 = \frac{\omega_1}{p}, \quad \dotsc, \quad \omega_\varepsilon = \frac{\omega_{\varepsilon-1}}{p}
\]
and if $\varphi$ denotes a rational function completely determined by the
class $l$:
\[
\begin{aligned}
M &= \left(\frac{\eta(\omega_1)}{\eta(\omega)}\right)^{24} &&= \varphi(k,\, lk), \\
(2)\quad M_1 &= \left(\frac{\eta(\omega_2)}{\eta(\omega_1)}\right)^{24} &&= \varphi(lk,\, l^2k), \\
&\dotsc\dotsc\dotsc\dotsc\dotsc\dotsc\dotsc\dotsc\dotsc\dotsc\dotsc\dotsc \\
M_{\varepsilon-1} &= \left(\frac{\eta(\omega_\varepsilon)}{\eta(\omega_{\varepsilon-1})}\right)^{24} &&= \varphi(l^{\varepsilon-1}k,\, k).
\end{aligned}
\]
By multiplication of all these equations there follows:
\weq{3}{\left(\frac{\eta(\omega_\varepsilon)}{\eta(\omega)}\right)^{24} = \varphi(k,\, lk)\,\varphi(lk,\, l^2k)\,\dotsc\,\varphi(l^{\varepsilon-1}k,\, k).}
But by § 120 the right-hand side of this equation can be represented as a
rational function $\Phi(k, lk)$ of $(k)$ and $(lk)$ (with rational numerical
coefficients).
On the other hand, since $l$ belongs to the exponent $\varepsilon$, $l^\varepsilon$ is equivalent to
the principal form, and so $l^\varepsilon k$ to $k$ and $\omega_\varepsilon$ to $\omega$; there is therefore
an equation:
\weq{4}{\omega_\varepsilon = \frac{\gamma + \delta\,\omega}{\alpha + \beta\,\omega}, \quad \alpha\delta - \beta\gamma = 1,}
and $\eta(\omega_\varepsilon) : \eta(\omega)$ is to be determined according to § 38. If, as
we have assumed, $\omega$ satisfies the equation
\weq{5}{a\,\omega^2 + b\,\omega + c\,p^\varepsilon = 0,}
then one obtains from (4), since $\omega_\varepsilon = \omega/p^\varepsilon$:
\[
\beta\,\omega^2 + (\alpha - p^\varepsilon\delta)\,\omega + \gamma\,p^\varepsilon = 0,
\]
which by comparison with (5) gives
\[
\beta = a\,x, \quad \alpha - p^\varepsilon\delta = b\,x, \quad \gamma = -c\,x,
\]
and if one sets $\alpha + p^\varepsilon\delta = y$:
\weq{6}{\alpha = \frac{y + b\,x}{2}, \quad \beta = a\,x,}
\[
p^\varepsilon\gamma = -c\,x, \quad p^\varepsilon\delta = \frac{y - b\,x}{2},
\]
%CONT>>
